\documentclass[a4paper,10pt]{article}
\usepackage[margin=1.5cm]{geometry}
\usepackage{multicol}
\usepackage{hyperref}
\usepackage{enumitem}
\usepackage{listings}



\setlength{\parindent}{0pt}

\begin{document}


\begin{center}
    {\Large \texttt{HTTP(S)}}\\
\end{center}

\tableofcontents
\newpage



\section{¿Qué es HTTP(S)?}


\textbf{HTTP} (Hypertext Transfer Protocol) es el protocolo encargado de la transferencia de infromación a través de archivos (XML, HTML, etc.).
HTTP define la sintaxis y la semántica que utilizan los elementos de software de la arquitectura web (clientes, servidores, proxies) para comunicarse.

\

Es un protocolo orientado a transacciones y sigue el esquema petición-respuesta entre un cliente y un servidor. El \textbf{cliente}  
realiza una petición enviando un mensaje, con cierto formato al servidor. El \textbf{servidor} le envía un mensaje de respuesta.


\section{Estructura de un mensaje HTTP}

Los mensajes HTTP son en texto plano, lo que lo hace más legible y fácil de depurar. Los mensajes tienen la siguiente estructura:

\begin{itemize}
    
    \item Línea inicial (termina con retorno de carro y un salto de línea) con:
    
    \begin{itemize}
        \item[*] Para las \textbf{peticiones}: la acción requerida por el servidor (método de petición) 
        seguido de la URL del recurso y la versión de HTTP que soporta el cliente.
        
        \item[*] Para \textbf{respuestas}: la versión del HTTP usado seguido del código de respuesta 
        (que indica qué ha pasado con la petición seguido de la URL del recurso) y de la frase asociada a dicho retorno.
    
    \end{itemize}
    
    \item Las \textbf{cabeceras} del mensaje que terminan con una línea en blanco. Son metadatos. Estas cabeceras le dan gran flexibilidad al protocolo.
    
    \item Cuerpo del mensaje (\textit{Opcional}). Su presencia depende de la línea anterior del mensaje y del tipo de recurso al que hace referencia la URL. 
    Típicamente tiene los datos que se intercambian cliente y servidor. 
    Por ejemplo, para una petición podría contener ciertos datos que se quieren enviar al servidor para que los procese. Para una respuesta podría incluir los datos que el cliente ha solicitado.

\end{itemize}




\section{Métodos}

HTTP define una \textbf{serie predefinida de métodos} de petición para indicar la acción que se desea realizar sobre el \textbf{recurso} identificado. 
Que este recurso sea preexistente (datos ya almacenados) o se genere de forma dinámica depende de la implementación del servidor. 
A menudo, el recurso corresponde a un archivo o a la salida de un ejecutable que reside en el servidor.

\

Un método puede ser catalogado como \textbf{seguro, idempotente o cacheable}:

\begin{itemize}

    \item Método Seguro:  Si una petición con ese método no tiene un efecto intencional en el servidor (\textbf{solo lectura}).
    
    \item Método idempotente: Si múltiples peticiones con ese método tienen el mismo efecto que una sola petición.
    Los métodos seguros son trivialmente idempotentes pues no deberían representar actualizaciones para el servidor.
    
    \item Método Cacheable: Si se pueden almacenar sus resultados para su reutilización.

\end{itemize}

\begin{tabular}{| c | c | c | c | c | c |}
    \hline
    Método & Petición con Carga Util & Respuesta con Carga Util & Safe & Idempotente & Cacheable \\
    \hline
    GET & Opcional	& Sí & Sí & Sí & Sí \\ 
    \hline
    HEAD & Opcional	& No & Sí&	Sí&	Sí\\
    \hline
    POST & Sí &Sí &No	&No	&Sí\\
    \hline
    PUT & Sí &Sí &No	&Sí	&No\\
    \hline
    DELETE &Opcional	&Sí	&No	&Sí	&No\\
    \hline
    CONNECT &Opcional	&Sí &No	&No	&No \\
    \hline
    OPTIONS& Opcional&	Sí&	Sí&	Sí	&No \\
    \hline
    TRACE &  No	&Sí	&Sí	&Sí	&No \\
    \hline
    PATCH & Sí	&Sí	&No	&No	&No\\
    \hline
    QUERY & & &Sí & Sí & Sí\\
    \hline
\end{tabular}

\

\subsection{GET}
El método \textbf{GET} solicita que el recurso de destino transfiera una representación de su estado. 
Las solicitudes que usan GET solo deben recuperar datos y no deben tener ningún otro efecto. \textit{Traeme el recurso}.

\subsection{HEAD}
El método \textbf{HEAD} solicita que el recurso de destino transfiera una representación de su estado, al igual que el método GET, pero sin los datos de representación contenidos en el cuerpo de la respuesta.
Esto es útil para poder recuperar los metadatos de las cabeceras de la respuesta, sin tener que transportar el contenido completo.
\textit{Traeme SOLO la etiqueta del recurso}.

\subsection{POST}
El método \textbf{POST} solicita que el recurso de destino procese la representación contenida en la petición de acuerdo con la semántica del propio recurso. 
\textbf{A nivel semántico está orientado a crear un nuevo recurso}.
\textit{Crea el recurso que con los  datos que te voy a pasar}.

\subsection{PUT}
El método \textbf{PUT} envía datos al servidor, pero a diferencia del método POST la URI de la línea de petición no hace referencia al recurso que los procesará, sino que identifica a los propios datos. 
PUT está \textbf{más orientado a la actualización} de los mismos (aunque también podría crearlos).
\textit{Actualiza (o crea) el recurso con los siguientes datos}.

\subsection{DELETE}
Borra el recurso.

\subsection{CONNECT}
Es una solicitud que establece un \textbf{Tunel TCP/IP} medinate un proxy HTTP, con el fín de facilitar comunicación mediante \textbf{HTTPS}.
\textit{Solicito establecer conexión contigo}. 

\subsection{OPTIONS}
Se usa para \textbf{mostrar las opciones presentes durante una conexión para un URL especifico}, con el fín de que el cliente sea capaz de determinar qué métodos y/o cabeceras HTTP están permitidas
durante una conexión.
\textit{¿Qué solicitudes te puedo hacer respecto a este recurso?}

\subsection{TRACE}
El método \textbf{TRACE} Solicita al servidor que introduzca en la respuesta todos los datos que reciba en el mensaje de petición. 
Se utiliza con fines de depuración y diagnóstico ya que el cliente puede ver lo que llega al servidor y de esta forma ver todo lo que añaden al mensaje los servidores intermedios 
\textit{Dime TODO lo que los servidores intermedios añaden a mi mensaje}.

\subsection{PATCH}
Su función es la misma que PUT, el cual \textbf{sobrescribe completamente un recurso}. Se utiliza para actualizar, de manera parcial una o varias partes. Está orientado también para el uso con proxy.
\textit{Actualiza el recurso}.

\subsection{QUERY}
Similar a GET, \textbf{QUERY} solicita un recurso con la diferencia de poder solicitarlo mediante estructuras de consulta (queries) como los son \textbf{JSON o SQL}.
Esto elimina la limitación del tamaño de URLs o la no idempotencia de la semantica de POST.
\textit{Busca el recurso utilizando este JSON o esta consulta SQL}.



\section{Ejemplo de Dialogo}

\begin{itemize}
    \item Cliente: 
        \begin{lstlisting}
            GET /index.html HTTP/1.1
            Host: www.example.com
            Referer: www.google.com
            User-Agent: Mozilla/5.0 (X11; Linux x86_64; rv:45.0) Gecko/20100101 Firefox/45.0
            Connection: keep-alive
            [Linea en blanco]
        \end{lstlisting}

    \item Servidor:
    \begin{lstlisting}
            HTTP/1.1 200 OK
            Date: Fri, 31 Dec 2003 23:59:59 GMT
            Content-Type: text/html
            Content-Length: 1221

            <html lang="eo">
            <head>
            <meta charset="utf-8">
            <title>Titulo del sitio</title>
            </head>
            <body>
            <h1>Pigina principal de tuHost</h1>
            (Contenido)
            .
            .
            .
            </body>
            </html>
        \end{lstlisting}
\end{itemize}


Para ver un \textbf{ejemplo de dialogo de cada uno de los métodos HTTP}, desde un shell linux:
\begin{lstlisting}
    curl -X [Metodo] https://example.com/resource (-d '{"key":"value"})'
\end{lstlisting}



\section{Status Codes}
Un \textbf{STATUS CODE} es un número que indica que ha pasado con la petición. 
El resto del contenido de la respuesta dependerá del valor de este código.
\

El número de los códigos están elegidos de tal forma que según si pertenece a una centena u otra se pueda identificar 
el tipo de respuesta que ha dado el servidor:

\begin{itemize}
    \item Formtado \textbf{1XX: LOGS}. Respuestas informativas. Indica que la petición ha sido recibida y se está procesando.
    \item Formato \textbf{2XX: Respuestas correctas}. Indica que la petición ha sido procesada correctamente.
    \item Formato \textbf{3XX: Respuestas de redirección}. Indica que el cliente necesita realizar más acciones para finalizar la petición. 
    |   \textit{Lo que buscas ya no está aquí}.
    \item Formato \textbf{4XX: Errores del lado del cliente}. Indica que ha habido un error al procesar le petición del cliente pues algo fue mal de su lado.
    \item Formato \textbf{5XX: Errores del lado del servidor}. La petición llegó, pero no se pudo procesar debido a un fallo del lado del servidor. 
\end{itemize}

\subsection{Status Codes Comunes}

\begin{tabular}{| c | c | c |}
    \hline
    Código & Significado & Descripción \\
    \hline
    200 & SUCCESS &  Petición Exitosa\\
    \hline
    301 & REDIRECT & El cliente será redirigido a otro dominio \\
    \hline
    400 & BAD REQUEST & \\
    \hline
    401 & UNAUTHORIZED & El cliente no ha sido autenticado \\
    \hline
    403 & FORBIDDEN & El cliente sí fue autenticado, más no tiene acceso al recurso. \\
    \hline
    404 & NOT FOUND & El recurso solicitado no pudo ser encontrado. \\
    \hline
    500 & INTERNAL SERVER ERROR & Problemas del lado del servidor. \\
    \hline
    503 & SERVICE UNAVAILABLE & Servidor no disponible. \\
    \hline
\end{tabular}

\end{document}
